\section{Multiphase Electric Machine Requirements}

The SPB converter comprises series-connected submodules, each containing a conventional two-level, three-phase converter. Its ac terminals can therefore be connected to a multistar load, including a multiphase machine \cite{Jin2017}. For the demonstrated implementation, the preferred machine architecture was a multiphase winding system that allowed the submodules to control separate three-phase winding sets. This observation establishes compatibility for the reported machine, but does not show that all SPB implementations require independently controllable winding groups or exclude alternative winding connections.

The experimentally demonstrated machine was a four-pole permanent-magnet-assisted synchronous-reluctance machine with 12 available windings. These windings were arranged as four three-phase sets, and each of the four submodules controlled one set \cite{Jin2017}. Here, ``independent'' denotes independent converter control of the winding-set currents; it does not establish galvanic isolation or electromagnetic independence. The evidence does not quantify mutual inductance, spatial coupling, or asymmetry among the sets. Thus, the experiment demonstrates a particular converter--machine mapping rather than a universal requirement that the number of winding groups equal the number of converter modules.

The reported arrangement is consistent with modular electrical control because separate three-phase terminal sets are assigned to separate submodules. However, the available evidence does not establish the dc-side implications of this arrangement beyond the series-connected SPB topology, nor does it establish insulation, leakage-current, grounding, neutral-point, or packaging requirements. These properties depend on operating voltage, power level, grounding arrangement, and implementation conditions. % ADDITIONAL LITERATURE REQUIRED

The demonstrated current waveforms exhibited a $30^{\circ}$ electrical phase displacement attributed to the symmetrical winding configuration: the A1 and A3 groups were in phase, as were A2 and A4, while these two pairs were displaced by $30^{\circ}$ electrical \cite{Jin2017}. This result was obtained with a $15^{\circ}$ mechanical phase displacement, a $4$ Hz electrical speed, and a $200$ V dc bus \cite{Jin2017}. These conditions are essential for interpreting the measured phase relationship. The experiment verifies the specified current displacement under those operating conditions; it does not by itself quantify the effects of displacement on torque production, air-gap-field distribution, torque ripple, or efficiency. Those effects require a validated machine model or additional measurements. % ADDITIONAL LITERATURE REQUIRED

For the reported implementation, four submodules were mapped to four controllable three-phase winding sets and hence to 12 available windings \cite{Jin2017}. This mapping should not be generalized to all SPB-fed machines because alternative winding connections or converter--machine mappings are not evaluated by the supplied evidence. More generally, the phase number and winding arrangement introduce design questions involving terminal count, winding placement, insulation, interconnection, control effort, and symmetry. Quantitative comparisons with conventional three-phase traction machines for winding complexity, losses, power density, efficiency, and manufacturing burden require operating-condition-matched evidence. % ADDITIONAL LITERATURE REQUIRED

The demonstrated phase displacement also indicates that the winding sets cannot be analyzed solely as unrelated three-phase machines. Their spatial relationship must be included when machine currents are coordinated. Nevertheless, the supplied evidence does not provide a complete multiphase transformation, harmonic-subspace decomposition, sequence-component analysis, or separation of torque-producing and non-torque-producing components. Formal treatment of these quantities requires additional machine-theory or validated modeling literature. % ADDITIONAL LITERATURE REQUIRED

Similarly, independent current control at the converter level does not establish electromagnetic decoupling. Mutual inductance, spatial harmonics, winding asymmetry, and control interactions may influence the aggregate machine field and current regulation, but their magnitudes and effects on torque ripple, losses, and efficiency are not quantified in the reported experiment. The experiment therefore demonstrates controllable current waveforms for the four winding sets, while the adequacy of local control alone and the need for supervisory coordination remain unresolved design questions. % ADDITIONAL LITERATURE REQUIRED

Multiple controllable winding groups may permit conditional post-fault operation, but no fault-tolerance advantage is experimentally established by the supplied evidence. The cited experiment does not demonstrate converter-module faults, winding open circuits, winding short circuits, or post-fault torque production. Any such capability would depend on winding geometry, electrical isolation, converter voltage and current limits, thermal limits, and the reference-redistribution strategy. Quantitative fault derating and post-fault torque, voltage, current, and thermal limits require additional evidence. % ADDITIONAL LITERATURE REQUIRED

Accordingly, the demonstrated machine requirement is narrower than a general fault-tolerant or high-performance multiphase-drive prescription: the reported SPB implementation used a four-pole, 12-winding machine arranged as four controllable three-phase sets, with a measured $30^{\circ}$ electrical current displacement under specified test conditions \cite{Jin2017}. The experiment verifies this converter--machine compatibility. It does not, by itself, establish universal winding requirements, electromagnetic decoupling, fault tolerance, or improvements in reliability, power density, efficiency, or maintainability. Resolving those issues requires comparative machine studies, validated multiphase models, and faulted-operation experiments. % ADDITIONAL LITERATURE REQUIRED